\section{Программная реализация алгоритма DeepSORT с использованием расширенного фильтра Калмана}

Программа реализует алгоритм DeepSort, который предназначен для отслеживания объектов в видео или потоке изображений с использованием метода глубокого обучения для извлечения признаков и метрики ближайшего соседа для сопоставления между кадрами~\cite{Sun2006}. Основной задачей алгоритма является связывание объектов между последовательными кадрами, обеспечивая корректное и стабильное их отслеживание в реальном времени.

\subsection{Используемые программные компоненты}

Класс DeepSort, предназначенный для отслеживания объектов в реальном времени, включает в себя следующие нестандартные модули и классы:

\begin{itemize}
    \item Модуль для извлечения признаков объектов из их изображений с использованием заранее обученной нейросети;
    \item Модуль, реализующий метрику косинусного расстояния для сопоставления объектов по визуальным признакам;
    \item Класс, который представляет детекцию объекта, включая его координаты, уровень уверенности и признаки;
    \item Модуль управлениями состояниями треков (инициализация новых треков, предсказание их положения и обновление). Данный модуль представляет наибольший интерес, так как в нем мы реализуем расширенный фильтр Калмана, благодаря которому алгоритм по отслеживанию объектов получает возможность работать с нелинейными состояниями окружающих ТС
\end{itemize}

\subsection{Алгоритм работы программы DeepSORT}

Алгоритм работы программы DeepSORT начинается с инициализации, где задаются ключевые параметры, необходимые для корректного отслеживания объектов. В их числе: путь к модели для извлечения признаков максимально допустимое косинусное расстояние между признаками объектов, порог уверенности детекции, ниже которого объекты игнорируются, порог для подавления перекрывающихся объектов, количество кадров объект может отсутствовать перед удалением трека, лимит на количество сохраненных признаков для сопоставления и флаг использования GPU для ускорения работы.

После инициализации программа переходит к этапу обновления трекера. При каждом новом кадре извлекаются признаки объектов из изображений с помощью нейросети, преобразуются координаты объектов из формата центра и размеров  в формат верхнего левого угла, а затем генерируются объекты детекций, которые включают в себя координаты, уровень уверенности и признаки каждого объекта.

Затем трекер обновляется следующим образом: сначала выполняется предсказание новых положений активных треков на основе предыдущих данных, после чего выполняется обновление треков с учетом новых детекций. На завершающем этапе формируются выходные данные, включающие координаты объектов, идентификаторы треков и идентификаторы объектов. Эти данные собираются и возвращаются для дальнейшего использования.

Ключевые методы программы обеспечивают работу с координатами и признаками объектов. Присутствуют следующие методы: преобразование координат объектов из формата центра и размеров в формат верхнего левого угла и размеров; извлечение визуальных признаков из изображений объектов с использованием обученной модели; увеличение возраста треков в тех случаях, когда они не были обновлены в текущем кадре, что позволяет управлять удалением устаревших треков; методы преобразования координат между форматами верхнего угла и пар углов ограничивающего прямоугольника.

\subsection{Класс управления состояниями треков}

Алгоритм работы класса начинается с инициализации, где задаются ключевые параметры, необходимые для отслеживания объектов и управления их состояниями. В их числе: метрика расстояния для сопоставления объектов, максимальная допустимая IOU-дистанция при сопоставлении, максимальное количество кадров, в которых объект может отсутствовать перед удалением трека и количество кадров, необходимых для подтверждения нового трека.

После инициализации создается расширенный фильтр Калмана, который используется для предсказания и обновления положений объектов.

На каждом новом кадре вызывается метод предсказания, который обновляет предполагаемые позиции объектов с использованием расширенного фильтра Калмана. Далее увеличивается возраст каждого трека и помечается как пропущенный, если на текущем кадре не удалось найти соответствие.

Основное обновление треков выполняется следующим образом: сначала выполняется сопоставление текущих детекций с треками, который делит треки на подтвержденные и неподтвержденные, а затем использует комбинацию признаков и IOU для поиска соответствий. Треки, которым удалось найти соответствия, обновляются с помощью расширенного фильтра Калмана, пропущенные треки помечаются, а для новых детекций создаются новые треки. В конце старые или устаревшие треки удаляются.

Далее метод сопоставления делит процесс сопоставления на два этапа: сначала выполняется сопоставление подтвержденных треков с помощью метрики признаков, а затем оставшиеся треки сопоставляются по IOU. Это гарантирует, что треки остаются активными даже в случае временного отсутствия детекции.

Ключевые методы модуля:
\begin{itemize}
    \item Предсказывание нового положения объектов на основе  расширенного фильтра Калмана;
    \item Увеличение возраста треков и пометка пропущенных.
    \item Обновление треков на основе текущих детекций.
    \item Сопоставление детекций с треками.
    \item Создание нового трека для объекта, если он не нашел соответствия среди существующих.
\end{itemize}

